\documentclass[11pt,a4paper]{article}
\usepackage[utf8]{inputenc}
\usepackage[T1]{fontenc}
\usepackage{lmodern}
\usepackage[italian]{babel}
\usepackage{geometry}
\usepackage{url}
\usepackage{hyperref}
\usepackage{parskip}
\usepackage{csquotes}
\usepackage{graphicx}
\usepackage[backend=biber,style=authoryear,sorting=ynt]{biblatex}

\addbibresource{pubblicazioni.bib}

\geometry{margin=2.5cm}

\hypersetup{
    colorlinks=true,
    linkcolor=black,
    urlcolor=blue,
    citecolor=black
}

\title{Relazione delle attività di ricerca -- Terzo anno di dottorato in Patrimonio culturale nell'ecosistema digitale}
\author{Settembre 2024 -- Settembre 2025}
\date{}

\begin{document}

\maketitle

\section{Riassunto}

Il terzo anno del dottorato in Patrimonio culturale nell'ecosistema digitale ha visto da un lato la messa in produzione di HERITRACE, un framework che consente a curatori di dati del patrimonio culturale di curare i dati in maniera agnostica rispetto al \textit{data model}, e dall'altro il miglioramento e incremento di OpenCitations Meta, un indice di metadati bibliografici. Parallelamente, è proseguito il lavoro di ricerca sull'inversione dei grafi prodotti tramite RML (RDF Mapping Language), sulla gestione di dati temporali, sulle metodologie di estrazione di dati semantici e sullo sviluppo di infrastrutture di supporto.

\section{Attività principali}

\subsection{Sviluppo di HERITRACE}

HERITRACE è stata messa in produzione all'indirizzo \url{https://paratext.ficlit.unibo.it/} nell'ottobre 2024 nell'ambito del progetto ParaText, un'iniziativa di ricerca PRIN 2022 che studia le relazioni tra testi letterari e paratesti esegetici nella tradizione manoscritta della poesia greca antica. Il framework funziona come editor di dati semantici per il \textit{Database Bibliografico ParaText} e permette agli studiosi di Filologia Classica di curare relazioni bibliografiche complesse, correggere errori, migliorare e aggiungere informazioni, garantendo la tracciabilità delle modifiche e la registrazione della \textit{provenance} dei dati.

L'architettura del sistema è basata su Docker e supporta qualunque database a grafo o \textit{quadstore}. Per gestire le modifiche concorrenti è stato implementato un meccanismo di blocco delle risorse che utilizza Redis. L'autenticazione avviene attraverso l'integrazione con ORCID, permettendo di associare ogni modifica all'identità accademica del ricercatore che l'ha effettuata.

Il configuratore del sistema può definire il modello tramite SHACL, che viene utilizzato per la generazione dei form e per la validazione. È anche possibile non definire alcun \textit{data model}: il sistema funziona comunque, ma non è disponibile la validazione automatica. Le regole di visualizzazione possono essere configurate per singola classe o \textit{shape}, garantendo flessibilità nella presentazione dei dati.

Il sistema di ricerca alterna automaticamente tra query SPARQL regex e l'indice testuale, supportando al momento quelli di OpenLink Virtuoso e di Blazegraph, in base alle caratteristiche della ricerca effettuata. Il sistema è in grado di gestire entità orfane e cancellazioni a cascata preservando l'integrità referenziale.

Sono stati eseguiti test di usabilità tra agosto e settembre 2025 che saranno la base per i prossimi miglioramenti dell'usabilità e dell'esperienza utente.

È stato inoltre sviluppato un sistema di documentazione basato su Starlight/MDX per fornire documentazione tecnica completa di HERITRACE, facilitando così l'adozione da parte della comunità scientifica e la manutenzione del framework.

\subsection{OpenCitations Meta}

L'indice OpenCitations Meta è stato oggetto di significative ottimizzazioni delle prestazioni. Il sistema è stato migrato per operare esclusivamente attraverso triplestore, sia per i dati che per la \textit{provenance}, superando così le limitazioni dell'approccio precedente basato su file.

È stato sviluppato un sistema di fusione per risolvere il problema dei duplicati presenti nell'indice. Il software è in grado di identificare e unificare entità duplicate utilizzando esclusivamente gli identificatori condivisi come criterio di matching.

L'infrastruttura è stata arricchita con documentazione completa dei flussi di lavoro, test automatizzati e un sistema di versionamento semantico che prevede rilasci automatizzati, elementi fondamentali per garantire la manutenibilità della pipeline di elaborazione dati.

È stata inoltre sviluppata una suite completa di utility per Virtuoso che comprende automazione del deployment basata su Docker, strumenti per il caricamento massivo di dati e utility per l'ottimizzazione del database, al fine di supportare operazioni su dataset di grandi dimensioni all'interno dell'infrastruttura OpenCitations.

\subsection{Ricerca sulla gestione di dati temporali}

L'integrazione della Time Agnostic Library in HERITRACE ha evidenziato alcuni colli di bottiglia nelle prestazioni, in particolare nell'identificazione di entità cancellate. La funzionalità Time Vault di HERITRACE, che permette di visualizzare le entità cancellate, necessita infatti della ricostruzione di istantanee storiche per ciascuna entità eliminata al fine di determinarne le informazioni di classe, processo che genera significative limitazioni prestazionali quando applicato su larga scala.

La valutazione delle prestazioni della Time Agnostic Library è stata ostacolata dalla scarsa disponibilità di benchmark adeguati. Attualmente esistono soltanto tre benchmark specifici per sistemi RDF temporali: EvoGen, BEAR e SPBv. EvoGen presenta una limitata adozione da parte della comunità scientifica, rendendo impossibili confronti significativi tra diversi sistemi. I dataset di BEAR non risultano più accessibili online, nonostante i tentativi di contattare gli autori originali. SPBv, infine, richiede lo sviluppo di componenti Java ad hoc per l'interfacciamento con il framework di benchmark HOBBIT.

Lo sviluppo del componente Java necessario per l'integrazione con HOBBIT si è rivelato particolarmente complesso a causa dei requisiti specifici della piattaforma relativi all'interazione tra componenti e alla serializzazione dei dati. La piattaforma HOBBIT presenta inoltre problemi di stabilità: benchmark precedentemente pubblicati non riescono più a essere eseguiti correttamente sulla versione attuale della piattaforma, evidenziando problematiche sistemiche che hanno precluso una valutazione significativa delle prestazioni rispetto ad altri sistemi RDF temporali.

\subsection{RML e inversione di grafi di conoscenza}

La ricerca sull'inversione di grafi di conoscenza è stata svolta nell'ambito del programma di dottorato congiunto con KU Leuven sotto la supervisione di Anastasia Dimou. Questo studio affronta la problematica delle trasformazioni bidirezionali: mentre RML (RDF Mapping Language) permette di convertire database relazionali in RDF, il processo inverso – ovvero la conversione da RDF a strutture relazionali utilizzando le medesime definizioni di mappatura – non era mai stato approfondito.

Il sistema sviluppato riutilizza le mappature RML esistenti per effettuare trasformazioni inverse, garantendo coerenza tra conversioni dirette e inverse. Attraverso test approfonditi è stata verificata la copertura della maggior parte dei casi d'uso previsti dalle specifiche R2RML per database relazionali. Per la valutazione delle prestazioni è stato utilizzato il benchmark KROWN, originariamente sviluppato per processori RML come Morph-KGC. L'adattamento di questo benchmark ha permesso di quantificare tanto le caratteristiche di scalabilità del sistema quanto l'overhead computazionale delle trasformazioni bidirezionali rispetto all'elaborazione RML unidirezionale.

\subsection{Infrastruttura di supporto e strumenti}

È stato sviluppato l'estrattore SHACL SKG-IF per la generazione automatica di \textit{shape} SHACL a partire da ontologie annotate. Il sistema è in grado di analizzare i vincoli presenti nell'ontologia e di trasformarli automaticamente nelle corrispondenti \textit{shape} di validazione, eliminando così la necessità di creare manualmente tali \textit{shape}.

In collaborazione con PKP (Public Knowledge Project) e TIB (Technische Informationsbibliothek) è stato sviluppato un sistema di crowdsourcing per integrare i flussi di lavoro di invio delle citazioni all'interno di Open Journal Systems (OJS). Il sistema mira a colmare il divario di visibilità che caratterizza le 52.000 riviste OJS attive a livello mondiale, di cui il 79,9\% è localizzato nel Sud Globale ma risulta largamente assente dai principali indici di citazione.

Il flusso di lavoro implementato permette agli editori delle riviste di contribuire alle citazioni copiando i riferimenti bibliografici dagli articoli direttamente nei campi dedicati del dashboard OJS. Il sistema procede quindi all'analisi automatica dei riferimenti per estrarre metadati strutturati, offrendo la possibilità di correggere manualmente eventuali errori prima dell'invio finale. Le citazioni vengono quindi formattate come file CSV e inviate a GitHub sotto forma di issue strutturate. Il processo di validazione verifica la correttezza della struttura CSV, convalida gli identificatori e assicura la completezza dei metadati. Una volta al mese, i metadati vengono integrati in OpenCitations Meta e le citazioni in OpenCitations Index, mentre i metadati di \textit{provenance} e i dati vengono archiviati su Zenodo per la preservazione a lungo termine.

\section{Pubblicazioni}

\subsection{Articoli su rivista}

\fullcite{massariHERITRACEUserFriendlySemantic2025}

\fullcite{massariRepresentingProvenanceTrack2025}

\fullcite{barzaghiProposalFAIRManagement2024}

\fullcite{sebastianReproducibleWorkflowCreation2025}

\subsection{Software rilasciati}

\fullcite{arcangelomassariOpencitationsHeritraceV2332025}

\fullcite{arcangelomassariOpencitationsTimeagnosticlibrary5042025}

\fullcite{arcangelomassariOpencitationsOc_metaV2132025}

\subsection{Presentazioni}

\fullcite{massariMappingUnmappedCitation2025}

\fullcite{massariPresentationHERITRACEUserFriendly2025}


\vspace{2cm}

\noindent Ozzano dell'Emilia, 12 settembre 2025

\vspace{1cm}

\noindent \includegraphics[width=5cm]{firma.png}

\end{document}